% ============================================================================
% PROPOSED TEXT FOR AUTHOR REVIEW -- NOT AN AUTOMATIC PATCH
% Manuscript: Sensitivity of Fuzzy-PSO Inertia Control to Linguistic Granularity
%             and Rule-Base Structure (Mathematics)
% 14 October 2026 source PDF, 60 pages, used to preserve scientific scope.
% ============================================================================
% INTEGRATION INSTRUCTIONS FOR THE VS CODE AGENT:
% 1. These paragraphs are proposed replacements/adaptations for Introduction,
%    Related Work and Discussion. DO NOT blindly paste two \section{...} headings
%    into files that already define them. Preserve existing equation references,
%    H1--H5 organization, tables, and all numerical/statistical results.
% 2. Existing citations MUST use the keys in refs.bib shown below. The only two
%    MISSING references are Gacto2011 and Salimi2025; their BibTeX records are
%    supplied at the end as % comment-prefixed lines to be copied to refs.bib.
% 3. Existing BibTeX entries with inaccurate publication metadata are documented
%    at the end. Correct these EXISTING records, do not add duplicate keys.
% 4. The manuscript claims below are deliberately delimited to the studied
%    configurations and Classical F1--F23 benchmark. Do NOT introduce
%    confirmatory H3 inference, causal landscape explanations, or new results.
% 5. Incorporate paragraphs only after author approval and reconcile against the
%    final analysis tables generated by the VS Code agent.
% 6. Nothing herein authorizes changes to controller, experiments, CRediT,
%    data availability, raw results, appendices, or public repository.
%
% === PROPOSED SECTION 1: INTRODUCTION ===
\section{Introduction}
Particle swarm optimization (PSO) is a population-based method whose search behavior depends, among other design decisions, on the inertia weight used in the velocity update. An inertia schedule can regulate the relative influence of previous particle velocities throughout optimization, but a predefined schedule does not respond directly to changing swarm conditions. Fuzzy feedback provides one means of linking observed search indicators to parameter adjustments. Early work by Shi and Eberhart demonstrated fuzzy adaptation of PSO inertia weight, establishing a precedent for replacing a prescribed schedule with a linguistic control policy \cite{ShiEberhart2001}. Subsequent research has compared alternative inertia-adjustment strategies and proposed more elaborate fuzzy controllers for multiple PSO parameters \cite{Valdez2017,Xia2023,Roy2024}.

The effectiveness of a fuzzy controller depends not only on the variables selected as inputs and outputs, but also on how linguistic partitions and rule consequents are specified. The number of linguistic labels determines the resolution of the representation and, for a complete two-input rule table with equal granularity, the number of rules grows quadratically. The assignment of consequents then determines which nominal control responses are associated with combinations of linguistic states. These are distinct, but coupled, design choices. Research on linguistic fuzzy systems has long recognized the relevance of partition complexity and rule-base semantics to system interpretability \cite{Gacto2011,Castiello2019}. However, the relationship between increased representational detail and the performance of a fuzzy-controlled optimizer cannot be inferred from representational complexity alone.

Existing fuzzy-PSO studies address several complementary questions. Some compare complete parameter-adaptation strategies with standard or alternative PSO variants \cite{Valdez2017,Xia2023}; others optimize fuzzy feedback policies as part of algorithm design \cite{Roy2024}. Such work demonstrates the potential of fuzzy control, but evaluations that jointly change control policies, controlled parameters, or algorithmic components do not by themselves characterize the sensitivity of a fixed Mamdani inertia controller to linguistic granularity and consequent-pattern specification. The present study addresses this narrower question. It is not intended as a new PSO variant or an automated method for selecting fuzzy policies.

Specifically, we examine a two-input Mamdani fuzzy control system (FCS) that maps normalized swarm diversity and iteration progress to the inertia weight. Under a common inference mechanism, membership-function family, PSO implementation, and experimental budget, we vary four linguistic granularities ($n\in\{3,5,7,9\}$) and eight systematically generated rule-base consequent patterns ($R1$--$R8$), yielding 32 FCS configurations. The comparison includes a conventional PSO baseline and covers 62 function--dimension instances from the Classical F1--F23 benchmark. These design controls permit a focused comparison of the implemented configurations without claiming that the effects of the two factors have been independently or causally decomposed.

The contributions are fourfold. First, we provide a reproducible construction of linguistic partitions and rule-base families whose nominal consequent directions can be described and compared within one Mamdani architecture. Second, we document a factorial sensitivity study of the 32 configurations under matched benchmark and execution conditions. Third, we examine configuration rankings, function-level heterogeneity, inertia-weight trajectories, and comparisons with the PSO baseline, while distinguishing descriptive patterns from supported statistical conclusions. Fourth, we discuss how the observed variation across benchmark functions constrains the interpretation of a single globally preferred static configuration and motivates, without evaluating, subsequent research on configuration selection. Our conclusions apply to the implemented FCS and the benchmark considered; they do not establish a general interaction between the experimental factors or the superiority of an adaptive selector.

% === PROPOSED SECTION 2: RELATED WORK ===
\section{Related Work}
\subsection{Inertia-weight adaptation and fuzzy feedback in PSO}
The inertia weight regulates the contribution of a particle's previous velocity to its next movement. Time-dependent schedules offer a simple way to change that contribution during optimization, whereas state-dependent approaches seek to respond to the evolution of the search. Shi and Eberhart \cite{ShiEberhart2001} introduced fuzzy adaptation of the inertia weight using current best performance and the current inertia weight as inputs, with the change in inertia weight as the output. Their controller employed three linguistic labels and nine rules and was assessed on three benchmark functions. This seminal study established the feasibility of fuzzy inertia control but did not address the comparative behavior of several granularity levels and consequent-pattern families under a common controller architecture.

Valdez et al. \cite{Valdez2017} examined constant, random, linearly decreasing, nonlinearly decreasing, and fuzzy-adaptive inertia strategies, alongside fuzzy regulation of inertia and learning factors. Their proposed Mamdani controller used two inputs, three outputs, and a nine-rule base, with evaluation on mathematical functions and coevolutionary settings. The comparison primarily concerns inertia-adjustment approaches and parameter-control designs rather than a factorial variation of linguistic resolution and rule-consequent structure within one unchanged inference mechanism. It is therefore a close methodological precedent, but its central comparison differs from the present one.

More recent studies have extended fuzzy control to other components of PSO. Xia et al. \cite{Xia2023} used a multiple-input multiple-output controller to adjust learning-related parameters in a PSO variant combining fitness-based and novelty-based exemplars, with evaluation on CEC2017 problems. Salimi et al. \cite{Salimi2025} used diversity and iteration progress as fuzzy antecedents to adapt the acceleration coefficients $c_1$ and $c_2$ in a reachability-analysis application based on graph transformation systems. These contributions demonstrate the variety of control targets and optimization contexts for fuzzy feedback. They do not constitute direct experimental replications of a study that exclusively varies inertia-controller granularity and nominal consequent patterns.

\subsection{Linguistic granularity, rule structure, and interpretability}
In linguistic fuzzy systems, increasing the number of labels changes the granularity of fuzzy partitions and generally expands the potential rule base. Gacto et al. \cite{Gacto2011} organized interpretability criteria along two axes: complexity versus semantic interpretability and rule bases versus fuzzy partitions. This distinction is pertinent here because the number of labels and the allocation of rule consequents describe different properties of the FCS. Nevertheless, our benchmark evaluates optimization performance, not a validated measure of human interpretability. Accordingly, rule count and linguistic resolution are treated as structural descriptors, not direct interpretability scores.

Castiello et al. \cite{Castiello2019} proposed data-driven construction of interpretable fuzzy partitions with variable granularity for classification problems. Their work shows why selecting partition resolution can be a substantive modeling decision, but its classification-oriented objective and learned partitioning procedure differ from the fixed-family Mamdani controller evaluated here. Neither the evidence from classification nor conceptual arguments about finer partitions imply that increasing $n$ must improve PSO performance.

\subsection{Systematic fuzzy policy design and research positioning}
Roy et al. \cite{Roy2024} proposed a benchmark-driven framework for evolving fuzzy feedback policies and reported 127 fuzzy-PSO algorithms involving different combinations of controlled parameters. Their approach uses an optimized rule policy and Takagi--Sugeno--Kang-style fuzzy inference, targeting improved optimizer performance and reduced dependence on manually specified feedback rules. The number 127 refers to controlled-parameter combinations; it should not be equated with 127 levels or patterns of linguistic granularity comparable to the 32 configurations in this study. Their work is an important demonstration that fuzzy-policy design can be treated systematically, while the present contribution characterizes a fixed architectural setting instead of optimizing policies.

Across these studies, the comparison objectives differ: establishing feedback feasibility, improving complete PSO variants, optimizing fuzzy policies, or choosing interpretable partitions. We therefore focus on a narrower empirical question: under a shared two-input Mamdani inertia controller, how do predefined linguistic granularities and systematically generated consequent patterns relate to final optimization performance, ranking stability, and inertia dynamics over the evaluated functions? This distinction supports an experimental and explanatory contribution, not a claim that fuzzy-PSO configurations have never previously been studied systematically. It also does not imply a confirmed factorial interaction: the corresponding H3 evidence in this study is exploratory and descriptive.

% === PROPOSED REPLACEMENT/ENRICHMENT FOR SECTION 6: DISCUSSION ===
\section{Discussion}
\subsection{Rule-consequent structure and linguistic resolution}
The experiment indicates that the nominal direction and arrangement of rule consequents are associated with substantial differences in the rankings of the evaluated FCS configurations, whereas finer linguistic partitions do not yield a consistent improvement. This result is relevant to the distinction between fuzzy partition complexity and rule-base semantics discussed by Gacto et al. \cite{Gacto2011}. With a complete two-input table, increasing granularity from three to nine labels raises nominal rule cardinality from 9 to 81, but the additional representational capacity does not guarantee improved optimization outcomes under the tested conditions. This finding concerns optimization performance; it is not an empirical measurement of the interpretability--accuracy trade-off considered in linguistic modeling research \cite{Gacto2011,Castiello2019}.

The comparison complements prior fuzzy-PSO research rather than directly challenging its reported improvements. Shi and Eberhart \cite{ShiEberhart2001} provided an early fixed fuzzy policy for inertia adaptation, and Valdez et al. \cite{Valdez2017} compared several inertia-control strategies, including a hand-specified Mamdani rule base. Our experiment varies internal design choices within the same implemented inference architecture. It therefore makes visible differences that may be obscured when whole PSO variants or parameter-adjustment strategies are compared. At the same time, the fixed membership-function family and overlap structure limit any conclusion about the value of increased granularity in other fuzzy controllers.

The leading mean rank of FCS-3-R1 must be interpreted according to its implementation. R1 assigns identical consequents across diversity labels at a fixed progress label, but diversity still affects activation and aggregation under min--max Mamdani inference. Thus, the result does not isolate a causal advantage of diversity feedback, and its near-equal average inertia relative to the baseline does not make the two trajectories equivalent. These distinctions matter when moving from interpretable rule descriptions to claims about actual controller behavior.

\subsection{Structural symmetry and observed performance differences}
The mirror-rule construction provides a formal relationship between some effective control surfaces under the specified symmetric partition and inference conditions. The corresponding surfaces can be complementary around the inertia interval midpoint even when the associated optimizer outcomes differ. This contrast is useful because it separates a mathematical property of the controller mapping from an empirical property of the coupled PSO search process. Complementary output surfaces do not imply complementary or equivalent optimization performance. The proposition is a structural characterization within the adopted controller family, not a general theorem concerning all Mamdani systems or PSO dynamics.

\subsection{Configuration heterogeneity and relation to automated policy design}
The evaluated benchmark does not support treating the overall leading configuration as uniformly preferable for each function or class. Observed rank concordance between the scalable and fixed-dimensional subsets is weak, and the two evaluated subsets have different leading configurations. This heterogeneity is compatible with the broader motivation for state-dependent parameter control, but it does not identify which landscape property explains the rankings, prove statistical independence between subsets, or establish transferability beyond F1--F23.

Roy et al. \cite{Roy2024} approached the related challenge by optimizing fuzzy policies on benchmark problems, with generalization examined separately. Our results address a different prerequisite: documenting how predefined controller structures behave before developing or validating any online selection mechanism. An instance-wise oracle that selects the best policy after results are known cannot serve as evidence that a deployable selector will achieve the same benefit. Any future meta-selection approach must use only information available at decision time and must be tested on held-out problems to evaluate selection performance and avoid benchmark-informed leakage.

\subsection{Scope and limitations}
The present analysis uses uniform triangular partitions with shoulders, one fixed overlap convention, eight specified rule families, centroid defuzzification, and a common PSO implementation. Only the inertia weight is controlled. Accordingly, findings should not be extended without additional evidence to Gaussian or trapezoidal partitions, altered overlap degrees, alternative rule families, type-2 fuzzy systems, or PSO variants that adapt several parameters simultaneously. The Classical F1--F23 benchmark is not interchangeable with rotated or shifted CEC suites, and comparisons with modern competing optimizers were outside the stated experimental question.

The evidence concerning H3 remains descriptive. Neither the function-level rank dispersion nor the corresponding log-gap summaries show a consistent monotonic reduction across granularity levels; discrepant exploratory inferential procedures have not established a calibrated confirmatory interaction test for the dependence structure. The study consequently makes no claim that an interaction has been established or ruled out. Within these boundaries, the study provides a reproducible account of structural sensitivity and helps specify the questions that future adaptive fuzzy-configuration methods would need to answer.

% ============================================================================
% BIBLIOGRAPHY RECONCILIATION NOTES (FOR THE AGENT, NOT JOURNAL TEXT)
% Existing keys reused unchanged:
%   ShiEberhart2001, Valdez2017, Xia2023, Roy2024, Castiello2019
% New keys needed: Gacto2011, Salimi2025.
%
% IMPORTANT: fix metadata IN EXISTING refs.bib entries (DO NOT DUPLICATE):
%   Valdez2017: Applied Soft Computing 52 (2017), pp. 1070--1083
%               (current refs.bib incorrectly says 291--305).
%   Xia2023: IEEE Transactions on Fuzzy Systems 31(7) (2023), pp. 2464--2478
%            (current refs.bib incorrectly says 31(9), 3057--3071).
%   Castiello2019: Applied Soft Computing 74 (2019), pp. 567--582
%                  (current refs.bib incorrectly says 1053--1069).
%   Roy2024: same citation key is already used in manuscript; the published
%            reference is Roy, Beauthier and Mayer (2025), Evolutionary
%            Computation 33(2), pp. 279--308, DOI 10.1162/evco_a_00353.
%            KEEP the key Roy2024 to preserve \cite{} references; correct year
%            and publication fields IN ITS EXISTING bib entry.
%
% BibTeX: copy ONLY these two NEW RECORDS into refs.bib (without % prefixes):
% @article{Gacto2011,
%   author  = {Gacto, M. J. and Alcal{\'a}, R. and Herrera, F.},
%   title   = {Interpretability of Linguistic Fuzzy Rule-Based Systems: An Overview of Interpretability Measures},
%   journal = {Information Sciences},
%   year    = {2011},
%   volume  = {181},
%   number  = {20},
%   pages   = {4340--4360},
%   doi     = {10.1016/j.ins.2011.02.021}
% }
%
% @article{Salimi2025,
%   author  = {Salimi, Nahid and Soleimani, Seyfollah and Rafe, Vahid and Khodadad, Davood},
%   title   = {A Hybrid Approach for Reachability Analysis of Complex Software Systems Using Fuzzy Adaptive Particle Swarm Optimization Algorithm and Rule Composition},
%   journal = {Mathematical and Computational Applications},
%   year    = {2025},
%   volume  = {30},
%   number  = {3},
%   pages   = {65},
%   doi     = {10.3390/mca30030065}
% }
%
% No BibTeX entries were inserted directly into the user's refs.bib.
% Confirm metadata from original publisher PDFs before committing changes.
% ============================================================================
